\justifying
\NSFCBodyText

\subsubsection{与本项目相关的研究工作积累}
申请人长期从事时间序列分析、新能源功率预测及多尺度特征表示研究。围绕光伏功率中日周期变化与局部波动并存的特点，前期工作探索了小波分解与时序模型相结合的预测方法，积累了多尺度表示、非线性映射和时间依赖建模经验。这些积累与本项目的气象驱动功率建模直接相关，可用于建立可复现的预测基线，并分析引入气象信息后不同时间尺度上的误差变化。

本项目拟在上述积累上研究概率预测及爬坡风险。其新增任务包括：按实际预测时刻约束气象信息的可用范围，评估天气变化对预测区间覆盖与宽度的影响，以及利用相邻时刻的联合变化估计爬坡事件概率。当前尚无可核验的本项目专属先导实验，相关数据整理、概率校准与预警验证均作为拟开展工作安排。已有论文用于证明相关方法研究能力，本申请的效果仍需通过独立实验检验。

\subsubsection{已取得的研究工作成绩}
申请人作为通讯作者发表 DWT-Former 研究（\textit{Energy}, 2025, 139283；DOI：\nolinkurl{10.1016/j.energy.2025.139283}）\cite{sun2025dwt}，围绕小波多尺度特征与 Transformer 时间表示的融合开展光伏功率预测研究。该成果为本项目比较单尺度与多尺度预测、识别局部变化和周期成分的作用提供了相关方法积累。

申请人作为共同第一作者发表 WaveKAN 研究（\textit{Neurocomputing}, 2026, 133361；DOI：\nolinkurl{10.1016/j.neucom.2026.133361}）\cite{liu2026wavekan}，研究自适应小波分解与 Kolmogorov--Arnold 网络结合的长时序预测方法。该成果为处理非平稳时序和比较不同非线性预测模型提供了基础。两项成果均不能替代本项目在新疆场站上的实测检验，也不能作为已完成概率校准和爬坡预警的证据。

申请人主持自治区高校基本科研业务费科研项目“融合多源数据大模型驱动的光伏发电功率预测方法研究”（XJEDU2026P005，2026年1月至12月）。该项目与本申请具有光伏预测方向的关联。本申请拟将研究重点落实到多尺度气象影响、概率预测与爬坡风险，其任务和资助边界须依据\textbf{【待补 F-XJ-05：高校基本科研业务费限项核验及任务书逐项对照结果】}落实。既有论文不自动认定为该在研项目或本申请的阶段成果。

\subsubsection{可行性分析}
在理论上，光伏功率的周期变化、天气扰动与测量误差具有不同时间结构，分尺度比较可为识别气象影响和误差来源提供可检验的分析框架。概率预测可通过留出的时间段评价其覆盖率、区间宽度与事件概率误差；本项目不将总体平均覆盖率等同于各类天气均可靠，而是分别检验常见天气和快速变化条件下的偏差。对超出样本支持范围的天气情形，将报告不确定性与适用范围。

在技术上，前期小波与时序预测工作能够支持基线实现和模型对比。研究拟按时间先后划分训练、验证、初始校准和测试数据，确保插补、归一化和特征选择仅利用相应训练信息。小波等变换以预测时刻已观测窗口为输入，避免未来样本参与分解。分位数模型、简单统计基线与有限规模时序模型共同构成比较对象，通过消融实验判断性能变化来自气象信息、尺度表示还是校准步骤。

在人员与实施条件上，申请人具备相关方向的研究与研究生指导经验，已确认的 GPU 资源和软件工具链可以支撑计算实验。拟由申请人负责研究设计、实验规范和结果核查，参与人员与投入安排按工作条件部分列明的缺口落实。研究期限按2027年1月至2028年12月安排，先完成数据与方法基线，再开展校准、爬坡评价和原型整理。新疆场站实测数据及使用授权是地域验证的必要条件；业务 NWP 档案决定含气象预报输入的方案及日前扩展能否验证。仅有功率与历史观测时仍可开展短时主任务，其结论范围按实际输入条件限定。

\subsubsection{研究风险与应对措施}
\textbf{数据与进度风险。}新疆光伏实测数据或气象档案未按计划取得，将影响场景验证。拟在项目首季度核验授权、覆盖时段、时间戳和关键字段；若仍不能形成可用数据清单，则先在取得合法使用条件的公开数据上完成基线与评价流程，同时继续推进新疆数据获取。第一年度结束时若仍缺少足够的新疆数据，应明确缩小地域化验证范围，报告该目标未完成，不用公开数据或模拟功率替代新疆实测结论。

\textbf{概率预测技术风险。}天气分布变化和样本时间相关性可能使预测区间失准。若连续滚动评价窗口出现覆盖率偏离、区间显著变宽或爬坡概率校准恶化，拟比较滚动校准、按天气分组诊断与更简单的预测模型，分别报告不同情形下的结果。若复杂模型不能稳定优于基线，则保留较简洁模型及其适用范围，交付完整比较结果，不以训练集表现选择结论。

\textbf{罕见事件样本风险。}强爬坡或特定天气的独立事件过少，可能导致风险分级缺乏统计依据。若某一天气分层无法形成跨时间段重复评价，则回退总体历史模板，报告样本量与不确定性，不作该分层的定量优劣结论。非主要描述分组可按预定规则合并并明确标注；主事件定义、阈值和原始测试结果保持不变，避免为增加样本而调整判据。

\textbf{资源与合作风险。}业务数据授权、外部验证安排或共享算力的可用时间可能变化。拟优先采用已确认的计算资源，控制模型规模，以统一训练预算完成必要基线与消融。若合作验证不能落实，则保留可复现实验和离线原型，明确业务应用尚未验证；未经数据权利人许可，不公开原始场站数据。
